\section{Joint scan of deficits}\label{sec:v5-deficit}
A negative auxiliary amplitude describes a signal-shaped deficit, not a physical negative rate. This scan uses the existing fits.

\fig{0.97}{../figures/v5_combined_deficit_scan.pdf}{Profiled deficit scan. Top: observed signed local significance and the stress-background offset; shading marks negative fitted contributions. Middle: raw-root Gaussian reference and conditional Gaussian-response tails. Both assign one to nonnegative raw roots. Local triangles mark the $10^{-8}$ display floor. Bottom: two separate global orderings; triangles mark one-sided 95\% zero-count bounds. Direct error bars are central 95\% intervals.}{fig:deficit}

For $z=(r-a)/s$, the conditional local rule is $p^-=\Phi(z)$ when $r<0$, and one otherwise. The principal scan score is $T^-=\max_{m:r_m<0}(-z_m)$, with $-\infty$ for an empty set. The separate raw-depth statistic is $D^-=\max_m\max(0,-r_m)$.

The deepest raw deficit is at 72 MeV: $r=-3.490$, with uncalibrated local reference $2.41\times10^{-4}$. The largest stress-centered deficit is at 83 MeV: $r=-0.676$, $a=+7.707$, $s=0.983$. Its raw reference is 0.2495 but its conditional local value is $7.18\times10^{-18}$, mainly because of the stress offset.

The latter threshold has 0/200,000 GP and 0/1,000 direct exceedances, giving one-sided 95\% global bounds of $1.50\times10^{-5}$ and 0.003. Every simulated raw-depth maximum exceeds the observed deficit. These realizations are shared with the excess scan.

These illustrative tails follow the excess scan and do not adjust for choosing directions, methods or orderings. The unresolved tails establish neither a particle signal nor physical background validity.
